\documentclass[12pt,a4paper]{report}

% ==================== REQUIRED PACKAGES ====================
\usepackage[utf8]{inputenc}
\usepackage{times} % For Times New Roman font
\usepackage[top=1in, bottom=1in, right=1in, left=1.5in]{geometry} % Margins
\usepackage{setspace} % For line spacing
\usepackage{fancyhdr} % For custom headers/footers
\usepackage{graphicx} % For college logo insertion
\usepackage{titlesec} % For custom heading formats
\usepackage{tocloft} % For custom table of contents formatting
\usepackage{lipsum} % For dummy text placeholder (can be deleted)
\usepackage{listings} % For source code appendix
\usepackage{array} % For beautiful tables
\usepackage{caption} % For custom captions spacing
\usepackage{xcolor} % For professional syntax colors

% ==================== FORMATTING OPTIONS ====================
\doublespacing % Sets default line spacing to double (or custom 1.5)
\setstretch{1.5} % Strict 1.5 line spacing as per guideline
\setlength{\parskip}{6pt} % Paragraph spacing of 6pt

% Disable global hyphenation (preventing words from being cut at line ends)
\hyphenpenalty=10000
\exhyphenpenalty=10000

% Maintain strict, consistent margins without letting text overflow the right margin
\emergencystretch 3em

% Title formats: Do not underline, no colons
\titleformat{\chapter}[block]
  {\normalfont\fontsize{14}{16}\bfseries\filcenter\MakeUppercase}
  {\chapterpagename\ \thechapter}{1em}{}
\titleformat{\section}
  {\normalfont\fontsize{12}{14}\bfseries\MakeUppercase}
  {\thesection}{1em}{}
\titleformat{\subsection}
  {\normalfont\fontsize{12}{14}\bfseries} % Subsections in Camel Case
  {\thesubsection}{1em}{}

% Captions formatting (Size 10, Times New Roman)
\captionsetup{font={font=small,rm},labelfont=bf}

% Code highlighting config for Appendix
\lstset{
  language=Python,
  basicstyle=\ttfamily\footnotesize,
  breaklines=true,
  breakatwhitespace=true,
  showstringspaces=false,
  numbers=left,
  numberstyle=\tiny\color{gray},
  frame=tb,
  tabsize=4,
  columns=flexible,
  keepspaces=true,
  commentstyle=\color{green!50!black},
  keywordstyle=\color{blue}\bfseries,
  stringstyle=\color{orange!80!black},
  xleftmargin=15pt,
  xrightmargin=15pt
}

% ==================== HEADERS & FOOTERS ====================
\pagestyle{plain} % Set default style for preliminary pages to plain
\fancypagestyle{plain}{
  \fancyhf{}
  \fancyfoot[C]{\fontsize{10}{12}\selectfont \thepage} % Centered page number for preliminaries, no headers
  \renewcommand{\headrulewidth}{0pt}
  \renewcommand{\footrulewidth}{0pt}
}

% ==================== BEGIN DOCUMENT ====================
\begin{document}

% ==================== ROMAN PAGE NUMBERING FOR PRELIMINARIES ====================
\pagenumbering{roman}

% ==================== TITLE PAGE ====================
\begin{titlepage}
\begin{center}
    \setstretch{1.1}
    {\fontsize{14}{16}\selectfont \textbf{Industry Training}}\\[0.8cm]
    
    {\fontsize{16}{18}\selectfont \textbf{IMAGESHIELD - SECURE IMAGE \& STRING ENCRYPTION TOOL}}\\[0.8cm]
    
    {\fontsize{12}{14}\selectfont \textit{Submitted for partial fulfilment of the requirements for the award of the degree of}}\\[0.5cm]
    {\fontsize{14}{16}\selectfont \textbf{BACHELOR OF ENGINEERING}}\\[0.2cm]
    {\fontsize{12}{14}\selectfont \textit{in}}\\[0.2cm]
    {\fontsize{14}{16}\selectfont \textbf{COMPUTER SCIENCE AND ENGINEERING (IoT-CS-BCT)}}\\[0.6cm]
    
    {\fontsize{12}{14}\selectfont \textit{By}}\\[0.2cm]
    {\fontsize{12}{14}\selectfont \textbf{J. Reshmika (2451-24-749-022)}}\\[0.1cm]
    {\fontsize{12}{14}\selectfont \textbf{N. Meghana (2451-24-749-041)}}\\[0.1cm]
    {\fontsize{12}{14}\selectfont \textbf{P. Chaitanya Likhita (2451-24-749-045)}}\\[0.6cm]
    
    {\fontsize{12}{14}\selectfont \textit{Under the guidance of}}\\[0.2cm]
    {\fontsize{14}{16}\selectfont \textbf{M. Anupama}}\\[0.1cm]
    {\fontsize{12}{14}\selectfont Assoc. Prof, Department of CSE-ALLIED\\[0.5cm]}
    
    % College Logo - Replace 'mvsr_logo.png' with your logo file in Overleaf
    \includegraphics[width=2.5cm]{mvsrlogo.png}\\[0.5cm]
    
    {\fontsize{14}{16}\selectfont \textbf{MVSR ENGINEERING COLLEGE}}\\[0.1cm]
    {\fontsize{12}{14}\selectfont \textbf{(An Autonomous Institution)}}\\[0.1cm]
    {\fontsize{10}{12}\selectfont Department of Computer Science and Engineering\\
    (Affiliated to Osmania University \& Recognized by AICTE)\\
    Nadergul, Balapur Mandal, Hyderabad -- 501 510\\
    Academic Year: 2025-2026}
\end{center}
\end{titlepage}

% ==================== CERTIFICATE PAGE ====================
\newpage
\begin{center}
    \setstretch{1.2}
    {\fontsize{14}{16}\selectfont \textbf{MVSR ENGINEERING COLLEGE}}\\[0.2cm]
    {\fontsize{12}{14}\selectfont \textbf{(An Autonomous Institution)}}\\[0.2cm]
    {\fontsize{11}{13}\selectfont Department of Computer Science and Engineering\\
    (Affiliated to Osmania University \& Recognized\\
    by AICTE Nadergul, Balapur Mandal, Hyderabad -- 501 510}\\[1.0cm]
    
    % Logo placeholder
    \includegraphics[width=2.5cm]{mvsrlogo.png}\\[0.5cm] 
    
    {\fontsize{14}{16}\selectfont \textbf{CERTIFICATE}}\\[1.0cm]
\end{center}

\noindent This is to certify that the Industry Training work entitled \textbf{``IMAGESHIELD - SECURE IMAGE \& STRING ENCRYPTION TOOL''} is a bonafide work carried out by \textbf{J. Reshmika (Roll No: 2451-24-749-022)}, \textbf{N. Meghana (Roll No: 2451-24-749-041)}, and \textbf{P. Chaitanya Likhita (Roll No: 2451-24-749-045)} in partial fulfilment of the requirements for the award of degree of \textbf{Bachelor of Engineering} in \textbf{Computer Science and Engineering (IoT-CS-BCT)} from \textbf{MVSR Engineering College}, affiliated to \textbf{OSMANIA UNIVERSITY}, Hyderabad, during the Academic Year 2025-2026 under our guidance and supervision.\\[0.8cm]
\noindent The results embodied in this report have not been submitted to any other university or institute for the award of any degree or diploma.\\[2.0cm]

\noindent \textbf{Internal Guide} \hfill \textbf{Head of the Department}\\[0.5cm]
\noindent M. Anupama \hfill Dr. [HOD Name] \\
\noindent Assoc. Professor, Dept. of CSE \hfill Professor \& Head, Dept. of CSE-Allied

% ==================== DECLARATION PAGE ====================
\newpage
\begin{center}
    {\fontsize{14}{16}\selectfont \textbf{DECLARATION}}\\[1.5cm]
\end{center}

\noindent We hereby declare that the work reported in the present Industry Training entitled \textbf{``IMAGESHIELD - SECURE IMAGE \& STRING ENCRYPTION TOOL''} is a record of bonafide work done by us in the Department of Computer Science and Engineering Allied, M.V.S.R. Engineering College, Osmania University. The reports are based on the work done entirely by us and not copied from any other source.\\[0.5cm]
\noindent The results embodied in this report have not been submitted to any other University or Institute for the award of any degree or diploma to the best of our knowledge and belief.\\[2.5cm]

\noindent
\begin{minipage}[t]{0.4\textwidth}
    \flushleft
    \textbf{Date:} \\
    \textbf{Place: Hyderabad}
\end{minipage}
\hfill
\begin{minipage}[t]{0.5\textwidth}
    \flushleft
    \textbf{Student 1: J. Reshmika} (2451-24-749-022)\\[0.2cm]
    \textbf{Student 2: N. Meghana} (2451-24-749-041)\\[0.2cm]
    \textbf{Student 3: P. Chaitanya Likhita} (2451-24-749-045)
\end{minipage}

% ==================== ACKNOWLEDGEMENT PAGE ====================
\newpage
\begin{center}
    {\fontsize{14}{16}\selectfont \textbf{ACKNOWLEDGEMENT}}\\[1.0cm]
\end{center}

\noindent We would like to express our sincere gratitude and indebtedness to our project guide \textbf{ M. Anupama}, Associate Professor, Department of Computer Science and Engineering, for her valuable suggestions and interest throughout the course of this project.\\[0.5cm]
\noindent We are also thankful to our principal       \textbf{       Dr. Marthi Kameswara Rao} and \textbf{Dr. Rajesh H. Kulkarni}, Professor and Head, Department of Computer Science and Engineering-Allied, MVSR Engineering College, Hyderabad for providing excellent infrastructure and a nice atmosphere for completing this project successfully as a part of our B.E. Degree CSE \textbf{(IoT-CS-BCT)}.\\[0.5cm]
\noindent We convey our heartfelt thanks to the lab staff for allowing us to use the required equipment whenever needed.\\[0.5cm]
\noindent Finally, we would like to take this opportunity to thank our family for their support through the work. We sincerely acknowledge and thank all those who gave directly or indirectly their support in completion of this work.\\[1.5cm]

\noindent
\hfill
\begin{minipage}[t]{0.5\textwidth}
    \flushleft
    \textbf{J. Reshmika} (2451-24-749-022)\\[0.2cm]
    \textbf{N. Meghana} (2451-24-749-041)\\[0.2cm]
    \textbf{P. Chaitanya Likhita} (2451-24-749-045)
\end{minipage}

% ==================== VISION AND MISSION ====================
\newpage
\begin{center}
    {\fontsize{14}{16}\selectfont \textbf{VISION AND MISSION OF THE DEPARTMENT}}\\[1.0cm]
\end{center}

\noindent \textbf{VISION}\\[0.3cm]
To evolve as a premier academic and research hub aligned with the global standards in the emerging areas of Computer Science.\\[1.0cm]

\noindent \textbf{MISSION}\\[0.3cm]
\noindent \textbf{M1:} To impart strong fundamental and applied knowledge, producing technically competent engineers to address current industrial and societal challenges.\\[0.3cm]
\noindent \textbf{M2:} To enable students to develop feasible solutions with an awareness towards environment and sustainability.\\[0.3cm]
\noindent \textbf{M3:} To nurture a culture of research and innovation that contributes to cutting-edge advancements.\\[0.3cm]
\noindent \textbf{M4:} To foster holistic professional development with strong human and ethical values.

% ==================== COURSE CERTIFICATION ====================
\newpage
\begin{center}
    {\fontsize{14}{16}\selectfont \textbf{COURSE CERTIFICATION}}\\[1.0cm]
\end{center}

\noindent As a fundamental requirement of the BE CSE (IoT-CS-BCT) curriculum and to support the core competencies of the Industry Training program, all team members completed professional certification courses from globally recognized academies. These certifications validated software engineering skills and cryptographic protocols utilized during the development of \textbf{ImageShield}.\\[0.6cm]

\noindent \textbf{1. CISCO Networking Academy -- Python Essentials 1 and 2}\\[0.3cm]
\noindent This certification was successfully completed by all team members on \textbf{10th March 2026} and \textbf{25th March 2026}. The course validated Python core structures, multi-threading architectures, database management, and object-oriented implementations:

\begin{figure}[htbp]
    \centering
    \begin{minipage}{0.38\textwidth}
        \centering
        \includegraphics[width=\textwidth]{certificates/reshmika_python1.png}
        \caption*{J. Reshmika -- Essentials 1}
    \end{minipage}
    \hspace{1cm}
    \begin{minipage}{0.38\textwidth}
        \centering
        \includegraphics[width=\textwidth]{certificates/reshmika_python2.png}
        \caption*{J. Reshmika -- Essentials 2}
    \end{minipage}
    
    \vspace{0.15cm}
    
    \begin{minipage}{0.38\textwidth}
        \centering
        \includegraphics[width=\textwidth]{certificates/meghana_python1.png}
        \caption*{N. Meghana -- Essentials 1}
    \end{minipage}
    \hspace{1cm}
    \begin{minipage}{0.38\textwidth}
        \centering
        \includegraphics[width=\textwidth]{certificates/meghana_python2.png}
        \caption*{N. Meghana -- Essentials 2}
    \end{minipage}
    
    \vspace{0.15cm}
    
    \begin{minipage}{0.38\textwidth}
        \centering
        \includegraphics[width=\textwidth]{certificates/likhita_python1.png}
        \caption*{Chaitnya Likhita P -- Essentials 1}
    \end{minipage}
    \hspace{1cm}
    \begin{minipage}{0.38\textwidth}
        \centering
        \includegraphics[width=\textwidth]{certificates/likhita_python2.png}
        \caption*{Chaitnya Likhita P -- Essentials 2}
    \end{minipage}
    \caption{CISCO Networking Academy Python Essentials 1 and 2 Certificates}
    \label{fig:python_certificates}
\end{figure}

\noindent \textbf{2. Simplilearn -- Introduction to Cryptography for Beginners}\\[0.3cm]
\noindent Mastered foundational cryptographic blocks, PBKDF2 stretching iterations, and HMAC authentication mechanisms deployed in the Security Controller of the ImageShield application:

\begin{figure}[htbp]
    \centering
    \begin{minipage}{0.38\textwidth}
        \centering
        \includegraphics[width=\textwidth]{certificates/reshmika_crypto.png}
        \caption*{J. Reshmika -- Cryptography}
    \end{minipage}
    \hspace{1cm}
    \begin{minipage}{0.38\textwidth}
        \centering
        \includegraphics[width=\textwidth]{certificates/meghana_crypto.png}
        \caption*{N. Meghana -- Cryptography}
    \end{minipage}
    
    \vspace{0.2cm}
    
    \begin{minipage}{0.38\textwidth}
        \centering
        \includegraphics[width=\textwidth]{certificates/likhita_crypto.png}
        \caption*{Chaitnya Likhita P -- Cryptography}
    \end{minipage}
    \caption{Simplilearn Introduction to Cryptography for Beginners Certificates}
    \label{fig:crypto_certificates}
\end{figure}

% ==================== ABSTRACT PAGE ====================
\clearpage
\begin{center}
    {\fontsize{14}{16}\selectfont \textbf{ABSTRACT}}\\[1.0cm]
\end{center}

\noindent This report details the design and implementation of \textbf{ImageShield}, a highly secure desktop utility engineered to protect sensitive image and text data from unauthorized access and tampering. Written in Python, the system utilizes a high-contrast Obsidian dark theme Tkinter interface and deploys robust AES-128 cryptographic pipelines (via the Fernet algorithm) paired with custom HMAC-SHA256 data signing to assure data integrity. Keys are derived dynamically from user passwords through PBKDF2 stretching running over 100,000 SHA-256 iterations to defeat brute-force and dictionary attacks. \par
\noindent To defend against automated password cracking attempts, the application embeds a stateful Security Controller that enforces a strict 30-second system-wide lockout after three consecutive failed entries, featuring an active real-time countdown timer. The application also supports local database indexing via SQLite3 to maintain user archives without storing passwords or plain asset keys, alongside webcam stream capture utilities using OpenCV to securely grab camera frames for direct, local encryption. Additionally, the system integrates a secure standalone String Encrypter allowing portable text encryption, saving records as self-contained secure `.enc` JSON packages. The resulting application runs entirely offline, guaranteeing maximum user privacy.



% ==================== TABLE OF CONTENTS ====================
\newpage
\tableofcontents
\newpage
\listoftables
\newpage
\listoffigures

% ==================== ARABIC PAGE NUMBERING FOR MAIN CONTENT ====================
\newpage
\pagenumbering{arabic}
\setcounter{page}{1}

% Configure standard fancy headers and footers for main content chapters
\pagestyle{fancy}
\fancyhf{}
\fancyhead[R]{\fontsize{10}{12}\selectfont IMAGESHIELD - SECURE ENCRYPTION TOOL} % Header - Project Name (top right)
\fancyfoot[L]{\fontsize{10}{12}\selectfont Dept. of CSE-Allied, MVSREC(A)} % Footer - Dept on left
\fancyfoot[R]{\fontsize{10}{12}\selectfont \thepage} % Footer - Page number on right
\renewcommand{\headrulewidth}{0.4pt}
\renewcommand{\footrulewidth}{0.4pt}

% Redefine plain style for chapter start pages to have bottom-right numbering as well
\fancypagestyle{plain}{
  \fancyhf{}
  \fancyfoot[R]{\fontsize{10}{12}\selectfont \thepage}
  \fancyfoot[L]{\fontsize{10}{12}\selectfont Dept. of CSE-Allied, MVSREC(A)}
  \renewcommand{\headrulewidth}{0pt}
  \renewcommand{\footrulewidth}{0.4pt}
}

% ==================== CHAPTER 1: INTRODUCTION ====================
\chapter{Introduction}

\section{Problem Statement and Scope}
In the modern digital landscape, the local storage of sensitive personal or enterprise visual assets (such as private documents, medical records, financial statements, and investigative photos) introduces significant security vulnerabilities. Assets stored in plain-text formats (such as PNG or JPG) are susceptible to unauthorized extraction, offline data harvesting, or malware-driven modification. 
\par
The scope of the \textbf{ImageShield} project is to deliver a highly secure, responsive, and completely offline desktop platform for localized cryptographic storage. The tool implements password-based symmetric encryption to package sensitive images and arbitrary text blocks into custom secure archives, indexing them cleanly in a local relational database without storing plain password components or Master Keys on disk.

\section{Application Areas, Users and Challenges}
ImageShield addresses several crucial use cases:
\begin{itemize}
    \item \textbf{Personal Privacy}: Safeguarding personal identity documentation.
    \item \textbf{Healthcare Scans}: Protecting patient diagnostics (DICOM image exports).
    \item \textbf{Legal Evidence}: Ensuring chain of custody through absolute data integrity.
\end{itemize}
Key engineering challenges include preventing UI thread blocking during resource-heavy cryptographic stretching and securing text packages while preserving full portability.

% ==================== CHAPTER 2: SYSTEM REQUIREMENTS ====================
\chapter{System Requirements}

\section{Tools and Technologies Used in the Project}
The technology stack selected for the implementation of ImageShield is defined as follows:
\begin{itemize}
    \item \textbf{Language}: Python 3.8+ providing cross-platform runtime support.
    \item \textbf{Cryptographic Library}: \texttt{cryptography.io} implementing industry-standard Fernet (AES-128-CBC) and PBKDF2.
    \item \textbf{GUI Toolkit}: Tkinter utilizing standard theme engines for zero external UI runtime dependencies.
    \item \textbf{Image Processing}: Pillow (PIL) for image loading and thumbnail rendering.
    \item \textbf{Computer Vision}: OpenCV (\texttt{opencv-python}) for hardware camera interfacing.
    \item \textbf{Database}: SQLite3 providing a zero-configuration SQL engine.
\end{itemize}

\section{Hardware and Software Requirements}
The requirements for running the application are outlined below.

\begin{table}[h]
\centering
\caption{Software Requirements}
\begin{tabular}{|m{5cm}|m{8cm}|}
\hline
\textbf{Software Component} & \textbf{Specification} \\ \hline
Operating System & Microsoft Windows 10 / 11 \\ \hline
Python Version & Python 3.8 or higher \\ \hline
Cryptography Library & \texttt{cryptography} $\ge$ 41.0.0 \\ \hline
Imaging Toolkit & \texttt{Pillow} $\ge$ 10.0.0 \\ \hline
Computer Vision Library & \texttt{opencv-python} $\ge$ 4.10.0 \\ \hline
Database & SQLite3 (Built-in) \\ \hline
\end{tabular}
\end{table}

\begin{table}[h]
\centering
\caption{Hardware Requirements}
\begin{tabular}{|m{5cm}|m{8cm}|}
\hline
\textbf{Hardware Component} & \textbf{Minimum Specification} \\ \hline
Processor & Intel Core i3 / AMD Ryzen 3 or higher \\ \hline
Random Access Memory & 4 GB RAM minimum \\ \hline
Storage space & 100 MB free space (plus database storage) \\ \hline
Camera & Standard integrated or external USB Webcam \\ \hline
\end{tabular}
\end{table}

\section{I/O Specification}
\begin{itemize}
    \item \textbf{Inputs}: Image files (PNG, JPG, JPEG), plain string text, password (string), tags (string), and active webcam frame capture feeds.
    \item \textbf{Outputs}: Secure standalone `.enc` files (JSON packages containing ciphertext, base64 salt, and HMAC), database records, and recovered plaintext files.
\end{itemize}

% ==================== CHAPTER 3: SYSTEM ARCHITECTURE ====================
\chapter{System Architecture}

The system architecture of ImageShield is designed around a decoupled, three-tier desktop paradigm consisting of an interactive Frontend, a hardened Backend Orchestration Layer, and an offline localized Database. This architectural design ensures absolute isolation of the cryptographic pipeline from storage routines and user interactions.

\section{Architecture Diagram}
\begin{figure}[h]
\centering
\includegraphics[width=0.8\textwidth]{Architecture.png}
\caption{System Architecture of ImageShield}
\end{figure}

\par
As illustrated in Figure 3.1, the platform's execution flow and data processing are strictly compartmentalized into three primary architectural tiers:

\section{Frontend (Presentation Layer)}
The Frontend serves as the user-facing interface, managing interactive GUI components, visual feedback, and user event loops entirely offline. It comprises the following core components:
\begin{itemize}
    \item \textbf{User Interface}: The primary desktop window providing control panels for encryption and decryption workflows, password entries, and active result display blocks.
    \item \textbf{Gallery View}: A dynamic grid/list visualizer displaying all locally cataloged records, displaying real-time reconstructed thumbnails, original filenames, extraction dates, and metadata tags.
    \item \textbf{Image Display}: Comparative side-by-side preview panels capable of presenting the clean original image, the encrypted noisy representation, and the cleanly recovered decrypted outputs.
    \item \textbf{Lockout Timer}: A thread-safe real-time countdown widget that dynamically displays the remaining penalty lock duration when brute-force threshold is breached.
    \item \textbf{File Browser \& Status Messages}: Integrated system dialog file pickers alongside reactive status banners displaying success alerts, error warning popups, and attempt logs.
\end{itemize}

\section{Backend (Orchestration \& Security Layer)}
The Backend functions as the logical engine of the application, managing secure cryptographic processing, thread concurrency, input validation, and database operations routing. It contains:
\begin{itemize}
    \item \textbf{Image Input \& Output Handlers}: Utilities powered by the Pillow library that ingest raw images, validate file headers (such as JPEG/PNG signatures), convert files into raw byte arrays, and reconstruct recovered buffers safely to disk.
    \item \textbf{Key Management}: Handles key derivation using PBKDF2 stretching algorithms to generate 32-byte symmetric keys from human-readable passwords.
    \item \textbf{Encryption \& Decryption Engines}: Implement standard AES-128-CBC block ciphers with secure padding and HMAC-SHA256 signatures to render plaintext payloads as unreadable random noise and verify data authenticity prior to decryption.
    \item \textbf{Security Controller}: A stateful monitor that coordinates brute-force protection logic by tracking failed entries, managing active timer stamps, and resetting counters upon successful verification.
\end{itemize}

\section{Database (Offline Storage Layer)}
The Database layer provides localized, secure data archiving capability without transmitting assets over external networks. It contains:
\begin{itemize}
    \item \textbf{Encrypted Files Table}: The main relational database schema archiving the core encrypted BLOB data, filenames, original sizes, formats (JPEG/PNG), timestamps, and optional user tags.
    \item \textbf{Metadata Index}: A high-speed index facilitating sub-millisecond query responses for gallery loading and filtering.
    \item \textbf{Query Engine}: Handles localized SQL transactions including INSERT, SELECT, and DELETE operations, enabling filtering queries (e.g. search by filename keyword or sorting by date).
\end{itemize}

\section{Interface Data Flows}
Communication between layers is fully structured and secure:
\begin{itemize}
    \item \textbf{Frontend to Backend}: Frontend emits the file path, raw password keys, desired operation flags, and search keywords. Backend returns the processed byte streams, gallery array lists, lockout state triggers, and operational status updates.
    \item \textbf{Backend to Database}: Backend emits encrypted byte BLOBs, file tags, query commands, and timestamps. Database returns matching metadata records, specific archived BLOB arrays, and query transaction feedback.
\end{itemize}

\section{System Flow Chart}
The operational workflow of the ImageShield application is mapped systematically in Figure 3.2, tracing the sequential logic for both the encryption and decryption pathways. It details the execution routing from user image selection, KDF key derivation, and AES encryption to the stateful decryption validation flow that governs the password verification and brute-force lockout timers.

\begin{figure}[h]
    \centering
    \includegraphics[width=0.85\textwidth]{Flowchart.png}
    \caption{System Operational Flowchart (Encryption \& Decryption Processes)}
    \label{fig:system_flowchart}
\end{figure}

% ==================== CHAPTER 4: IMPLEMENTATION ====================
\chapter{Implementation}

This chapter details the modular software engineering design and developmental setup configured for the execution of ImageShield. It provides a comprehensive analysis of the local environment initialization alongside the technical functional descriptions of the core security, database, and cryptographic packages.

\section{Environmental Setup}
The environment is initialized by preparing a clean virtual environment, linking paths, and installing the specific dependency tree using the standard package installer:
\begin{verbatim}
pip install -r requirements.txt
\end{verbatim}

\section{Modules \& Their Description}
This section details the modular breakdown of the ImageShield software codebase, illustrating the separation of concerns between different processing layers. It defines the core responsibilities, data structures, and algorithmic routines associated with each critical system module.

The application contains six highly focused modules:
\begin{enumerate}
    \item \textbf{DatabaseModule}: Thread-safe SQLite tables managing metadata caching. \\
    It initializes local relational tables that catalog and manage the metadata index and encrypted files index completely offline. This establishes a highly organized local state representation that permits rapid retrieval without archiving plain-text master keys.
    
    \item \textbf{KeyManagementModule}: Handles password stretching using PBKDF2 over 100,000 iterations. \\
    It implements key derivation mechanisms that convert user passwords into high-entropy cryptographic keys. By utilizing high-grade iterations and randomized unique salts, the module successfully resists sophisticated offline dictionary and precomputed rainbow table attacks.
    
    \item \textbf{ImageInputOutputModule}: pillow normalizations and JSON package serializations. \\
    It handles the format-normalization, thumbnail scaling, and byte-array serialization routines for raw file buffers. It bundles text metadata and active payloads into compact, self-contained JSON packages that guarantee seamless portability across local execution targets.
    
    \item \textbf{EncryptionEngine}: AES-128-CBC block ciphers. \\
    It functions as the core symmetric encryption pipeline utilizing standard advanced block ciphers. This engine performs low-latency block-padding and byte-stream transformations to completely isolate and secure plain assets under high-strength, hardware-accelerated algorithms.
    
    \item \textbf{SecurityController}: Stateful lockout manager and HMAC-SHA256 signature generator. \\
    It oversees critical operational safety flows including brute-force protection and cryptographic data-signing validations. It dynamically tracks login failures to enforce time-based cooldown thresholds, preventing unauthorized automated intrusion attempts entirely.
    
    \item \textbf{LiveCameraModule}: CV2 capture loops. \\
    It interacts directly with the computer vision libraries to capture video streams and single frames from active local webcam hardware. The resulting raw frame arrays are immediately transferred to the secure processing pipeline, providing instant stream-to-disk encryption capabilities.
\end{enumerate}

% ==================== CHAPTER 5: TESTING AND RESULTS ====================
\chapter{Testing and Results}

\section{Test Case Scenarios}
Comprehensive verification matrix was executed to ensure system integrity.

\begin{table}[h]
\centering
\caption{Test Execution Matrix}
\begin{tabular}{|m{2cm}|m{4cm}|m{4cm}|m{3cm}|}
\hline
\textbf{Test ID} & \textbf{Scenario} & \textbf{Expected Result} & \textbf{Status} \\ \hline
TC-01 & Successful Image Encryption & Encrypted BLOB in DB and .enc package saved & PASS \\ \hline
TC-02 & Successful Image Decryption & Original image fully recovered and displayed & PASS \\ \hline
TC-03 & Incorrect Password Entry & Attempt recorded; remaining count displayed & PASS \\ \hline
TC-04 & Brute-Force Lockout & System lockout for 30s after 3 consecutive failures & PASS \\ \hline
TC-05 & String Encryption & Plain text encrypted and exported to TXT format .enc & PASS \\ \hline
TC-06 & Live Camera Capture & Live webcam stream loaded and preview window displayed & PASS \\ \hline
TC-07 & Live Video Frame Encryption & Webcam frame captured, successfully encrypted, and saved to DB & PASS \\ \hline
\end{tabular}
\end{table}

\clearpage
\section{Screen Shots of Test Cases}

\begin{figure}[h]
    \centering
    \includegraphics[width=0.75\textwidth]{tbp/en_success.png}
    \caption{TC-01: Successful Image Encryption with Status Modal Confirmation}
    \label{fig:tc01_encrypt}
\end{figure}

\begin{figure}[h]
    \centering
    \includegraphics[width=0.75\textwidth]{tbp/de_success.png}
    \caption{TC-02: Successful Image Decryption and Local Recovery Confirmation}
    \label{fig:tc02_decrypt_success}
\end{figure}

\begin{figure}[h]
    \centering
    \includegraphics[width=0.75\textwidth]{tbp/de_failed.png}
    \caption{TC-03: Failed Decryption Warning and Attempts Decrement Feedback}
    \label{fig:tc02_decrypt_fail}
\end{figure}

\begin{figure}[h]
    \centering
    \includegraphics[width=0.75\textwidth]{tbp/30sec.png}
    \caption{TC-04: Active 30-Second Brute-Force Lockout Countdown Timer}
    \label{fig:tc03_lockout}
\end{figure}

\begin{figure}[h]
    \centering
    \includegraphics[width=0.75\textwidth]{tbp/string.png}
    \caption{TC-05: Successful String Encryption and Plaintext Recovery Pipeline}
    \label{fig:tc04_string}
\end{figure}

\begin{figure}[h]
    \centering
    \includegraphics[width=0.75\textwidth]{tbp/camera_window.png}
    \caption{TC-06: Live Camera Capture Window Interface with Live Video Feed Preview}
    \label{fig:tc06_camera_window}
\end{figure}

\begin{figure}[h]
    \centering
    \includegraphics[width=0.75\textwidth]{tbp/camera_success.png}
    \caption{TC-07: Webcam Live Frame Encryption and Database Verification}
    \label{fig:tc07_camera_success}
\end{figure}

% ==================== CHAPTER 6: CONCLUSION ====================
\chapter{Conclusion \& Future Enhancements}

\section{Conclusion}
The implementation of the ImageShield platform demonstrates the feasibility and security of localized, offline cryptographic storage solutions for desktop systems. The project achieves several core conclusions:
\begin{itemize}
    \item \textbf{Highly Responsive Desktop Security}: Successfully designed, developed, and deployed a localized, entirely offline desktop utility that provides robust cryptographic storage for sensitive image files and portable string packages, mitigating risks associated with cloud storage and unauthorized external access.
    \item \textbf{Hardened Crypto Pipeline}: Engineered a strict cryptographic pipeline that leverages the industry-standard PBKDF2-SHA256 key derivation function with 100,000 stretching iterations to secure user passwords, directly integrated with AES-128 symmetric block ciphers via the Fernet algorithm to produce high-entropy encrypted payloads.
    \item \textbf{Active Intrusion Prevention}: Integrated a stateful Security Controller module that actively monitors decryption attempts to enforce a strict lockout mechanism and system-wide cooldown timer, effectively mitigating brute-force and automated dictionary attacks on protected archives.
    \item \textbf{Fluid User Experience}: Deployed an asynchronous concurrent architecture using background execution threads to isolate intensive cryptographic computations and database operations, ensuring that the Tkinter-based graphical user interface remains fully responsive without freezing or lagging during execution.
\end{itemize}

\section{Future Enhancements}
While the current version of ImageShield provides robust security safeguards, future development cycles will focus on expanding its features and compliance standards:
\begin{itemize}
    \item \textbf{Upgrade to AES-256}: Upgrading the symmetric key length and block cipher implementations from AES-128 to AES-256 ciphers, meeting the highest tier of government, industrial, and military-grade security compliance standards.
    \item \textbf{Hardware-Backed Biometrics}: Integrating operating system-level biometric authentication APIs, such as Windows Hello, to allow users to securely decrypt files using fingerprint recognition or facial scanners rather than relying solely on text-based passwords.
    \item \textbf{Secure Key Containers}: Developing protected local hardware key store integrations (such as TPM chips) and dedicated offline password manager vaults to safely store and manage derived keys without direct exposure to memory-space attacks.
    \item \textbf{Advanced Format Support}: Extending the ingestion pipeline and format-normalization modules to dynamically process, compress, encrypt, and preview complex data types including multi-page PDF files, high-resolution video streams, and structured ZIP archives.
\end{itemize}

% ==================== REFERENCES ====================
\begin{thebibliography}{99}
\bibitem{fips} NIST FIPS Pub 197: Federal Information Processing Standards (FIPS) Publication 197: Advanced Encryption Standard (AES). National Institute of Standards and Technology.
\bibitem{pbkdf2} IETF RFC 8018 (PKCS \#5): Password-Based Cryptography Specification (PBKDF2). Internet Engineering Task Force.
\bibitem{fernet} PyCa Cryptography: Fernet Symmetric Encryption and PBKDF2 Key Derivation Spec. Python Cryptography Authority (PyCa).
\bibitem{pillow} Pillow (PIL) Imaging Library: Pillow Image Processing, File Normalizations, and Buffer Handling API Reference.
\bibitem{opencv} OpenCV Library: OpenCV-Python Bindings: Hardware-Interfaced Webcam Capture Loops and Real-Time Frame Processing Specification.
\bibitem{sqlite} SQLite Consortium: SQLite Database Engine: Zero-Configuration Offline Local Relational Storage Models.
\end{thebibliography}

% ==================== APPENDIX ====================
\appendix
\chapter{Sample Source Code}

\section{PBKDF2 Key Derivation Method}
\begin{lstlisting}
class KeyManagementModule:
    """Handles password-based key derivation and key generation."""
    
    SALT_LENGTH = 16  % 16 bytes salt
    ITERATIONS = 100000  % PBKDF2 iterations
    
    @staticmethod
    def generate_salt():
        """Generate random salt for key derivation."""
        return os.urandom(KeyManagementModule.SALT_LENGTH)
    
    @staticmethod
    def derive_key_from_password(password, salt):
        """Derive 32-byte encryption key using PBKDF2."""
        kdf = PBKDF2HMAC(
            algorithm=hashes.SHA256(),
            length=32,
            salt=salt,
            iterations=KeyManagementModule.ITERATIONS,
            backend=default_backend()
        )
        key = kdf.derive(password.encode())
        return key
\end{lstlisting}

\section{HMAC-SHA256 Verification Method}
\begin{lstlisting}
class SecurityController:
    # ... (lockout methods)

    @staticmethod
    def generate_hmac(data, key):
        """Generate HMAC for data integrity verification."""
        hmac_obj = hmac.new(key, data, hashlib.sha256)
        return hmac_obj.hexdigest()
    
    @staticmethod
    def verify_hmac(data, key, hmac_tag):
        """Verify HMAC for data integrity."""
        calculated_hmac = SecurityController.generate_hmac(data, key)
        return hmac.compare_digest(calculated_hmac, hmac_tag)
\end{lstlisting}

\section{Symmetric Encryption \& Decryption Engines}
\begin{lstlisting}
class EncryptionEngine:
    """Handles image encryption using AES-128-CBC via Fernet."""
    
    @staticmethod
    def encrypt_image(image_bytes, encryption_key):
        """Encrypt image bytes using AES algorithm via Fernet."""
        try:
            # Generate Fernet key from derived key
            fernet_key = base64.urlsafe_b64encode(encryption_key)
            fernet = Fernet(fernet_key)
            
            # Encrypt image bytes
            encrypted_data = fernet.encrypt(image_bytes)
            return encrypted_data
        except Exception as e:
            print(f"Encryption error: {e}")
            return None


class DecryptionEngine:
    """Handles image decryption using AES algorithm via Fernet."""
    
    @staticmethod
    def decrypt_image(encrypted_data, encryption_key):
        """Decrypt image bytes using AES algorithm via Fernet."""
        try:
            # Generate Fernet key from derived key
            fernet_key = base64.urlsafe_b64encode(encryption_key)
            fernet = Fernet(fernet_key)
            
            # Decrypt image bytes
            decrypted_data = fernet.decrypt(encrypted_data)
            return decrypted_data
        except Exception as e:
            print(f"Decryption error: {e}")
            return None
\end{lstlisting}

\section{Stateful Lockout Brute-Force Protection}
\begin{lstlisting}
class SecurityController:
    """Manages password attempts and lockout mechanism."""
    
    FAILED_ATTEMPTS_LIMIT = 3
    LOCKOUT_DURATION = 30  # seconds
    
    def __init__(self):
        self.failed_attempts = 0
        self.lockout_until = None
    
    def is_locked_out(self):
        """Check if system is currently locked out."""
        if self.lockout_until is None:
            return False
        
        if datetime.now() >= self.lockout_until:
            self.lockout_until = None
            self.failed_attempts = 0
            return False
        return True
    
    def get_remaining_lockout_time(self):
        """Get remaining lockout time in seconds."""
        if not self.is_locked_out():
            return 0
        
        remaining = (self.lockout_until - datetime.now()).total_seconds()
        return max(0, int(remaining))
    
    def record_failed_attempt(self):
        """Record a failed decryption attempt."""
        self.failed_attempts += 1
        
        if self.failed_attempts >= self.FAILED_ATTEMPTS_LIMIT:
            self.lockout_until = (
                datetime.now() + timedelta(seconds=self.LOCKOUT_DURATION)
            )
            return (
                f"Too many failed attempts. "
                f"Locked out for {self.LOCKOUT_DURATION} seconds."
            )
        
        remaining = self.FAILED_ATTEMPTS_LIMIT - self.failed_attempts
        return f"Password incorrect. {remaining} attempts remaining."
    
    def reset_attempts(self):
        """Reset failed attempts on successful authentication."""
        self.failed_attempts = 0
        self.lockout_until = None
\end{lstlisting}

\section{Live Webcam Frame Capture Module}
\begin{lstlisting}
class LiveCameraWindow:
    """Live webcam preview window for capturing an image."""

    def _update_frame(self):
        """Refresh the live webcam frame."""
        if not self.active:
            return

        if not self.cap.isOpened():
            self.status_label.config(text="Camera disconnected.", fg="#d62828")
            self.window.after(200, self.close)
            return

        ok, frame = self.cap.read()
        if not ok:
            self.status_label.config(
                text="Unable to read from camera.", fg="#d62828"
            )
            self.window.after(60, self._update_frame)
            return

        self.current_frame = frame
        rgb_frame = cv2.cvtColor(frame, cv2.COLOR_BGR2RGB)
        image = Image.fromarray(rgb_frame)
        image.thumbnail((840, 520))
        self.current_photo = ImageTk.PhotoImage(image)
        self.video_label.config(image=self.current_photo, text="")
        self.video_label.image = self.current_photo

        if not self.captured_path:
            self.status_label.config(
                text="Live preview running. Click Capture Frame "
                     "when the image looks right.",
                fg="#b0b0b0"
            )

        self.window.after(30, self._update_frame)

    def capture_frame(self):
        """Save the current camera frame to a temporary PNG file."""
        if self.current_frame is None:
            messagebox.showerror("Camera Error", "No camera frame is available.")
            return

        rgb_frame = cv2.cvtColor(self.current_frame, cv2.COLOR_BGR2RGB)
        image = Image.fromarray(rgb_frame)

        timestamp = datetime.now().strftime('%Y%m%d_%H%M%S')
        temp_name = f"imageshield_camera_{timestamp}.png"
        temp_path = os.path.join(tempfile.gettempdir(), temp_name)
        image.save(temp_path, format="PNG")

        self.captured_path = temp_path
        self.status_label.config(
            text=f"Captured image saved to {temp_path}",
            fg="#06a77d"
        )
\end{lstlisting}

\section{Offline SQL Database Module}
\begin{lstlisting}
class DatabaseModule:
    """Handles all SQLite database operations for storing encrypted files."""
    
    def __init__(self, db_path='imageshield.db'):
        self.db_path = db_path
        self.conn = None
        self.lock = threading.Lock()
        self.init_database()
    
    def init_database(self):
        """Create database tables if they don't exist."""
        try:
            self.conn = sqlite3.connect(self.db_path, check_same_thread=False)
            cursor = self.conn.cursor()
            
            # Create encrypted files table
            cursor.execute('''
                CREATE TABLE IF NOT EXISTS encrypted_files (
                    id INTEGER PRIMARY KEY AUTOINCREMENT,
                    file_id TEXT UNIQUE NOT NULL,
                    original_filename TEXT NOT NULL,
                    display_name TEXT,
                    file_format TEXT NOT NULL,
                    encrypted_data BLOB NOT NULL,
                    file_size INTEGER NOT NULL,
                    encryption_date TIMESTAMP DEFAULT CURRENT_TIMESTAMP,
                    hmac_tag TEXT NOT NULL,
                    salt TEXT NOT NULL,
                    tags TEXT
                )
            ''')
            
            # Create metadata index table for fast queries
            cursor.execute('''
                CREATE TABLE IF NOT EXISTS metadata_index (
                    id INTEGER PRIMARY KEY AUTOINCREMENT,
                    file_id TEXT UNIQUE NOT NULL,
                    filename TEXT NOT NULL,
                    display_name TEXT,
                    format TEXT NOT NULL,
                    size INTEGER,
                    date TIMESTAMP,
                    FOREIGN KEY (file_id) REFERENCES encrypted_files(file_id)
                )
            ''')
            self.conn.commit()
        except sqlite3.Error as e:
            print(f"Database initialization error: {e}")
    
    def insert_encrypted_file(
        self, file_id, original_filename, file_format, encrypted_data,
        hmac_tag, salt, file_size, tags=None, display_name=None
    ):
        with self.lock:  % Thread-safe database access
            try:
                cursor = self.conn.cursor()
                cursor.execute('''
                    INSERT INTO encrypted_files 
                    (file_id, original_filename, display_name, file_format, 
                     encrypted_data, hmac_tag, salt, file_size, tags)
                    VALUES (?, ?, ?, ?, ?, ?, ?, ?, ?)
                ''', (
                    file_id, original_filename, display_name, file_format,
                    encrypted_data, hmac_tag, salt, file_size, tags
                ))
                
                # Update metadata index
                cursor.execute('''
                    INSERT INTO metadata_index 
                    (file_id, filename, display_name, format, size, date)
                    VALUES (?, ?, ?, ?, ?, ?)
                ''', (
                    file_id, original_filename, display_name, file_format,
                    file_size, datetime.now().isoformat()
                ))
                
                self.conn.commit()
                return True
            except sqlite3.Error as e:
                print(f"Error inserting encrypted file: {e}")
                return False
\end{lstlisting}

\section{Obsidian Neon Theme Custom UI Elements}
\begin{lstlisting}
class ImageShieldApp:
    """Main application window styles and button UI creation."""

    def setup_styles(self):
        """Set up modern, flat dark theme styles for Ttk widgets."""
        self.style = ttk.Style()
        self.style.theme_use('clam')
        
        # Frame styles
        self.style.configure(
            'TFrame', background=COLORS['bg_secondary'], borderwidth=0
        )
        self.style.configure(
            'AppFrame.TFrame', background=COLORS['bg_primary'], borderwidth=0
        )
        
        # Label styles
        self.style.configure(
            'TLabel', background=COLORS['bg_secondary'],
            foreground=COLORS['text_primary'], font=("Segoe UI", 10)
        )
        
        # Notebook (Tab container) style
        self.style.configure(
            'TNotebook', background=COLORS['bg_primary'], borderwidth=0
        )
        self.style.configure(
            'TNotebook.Tab',
            background=COLORS['bg_secondary'],
            foreground=COLORS['text_secondary'],
            font=("Segoe UI", 10, "bold"),
            padding=[20, 8]
        )
        self.style.map(
            'TNotebook.Tab',
            background=[('selected', COLORS['bg_primary'])],
            foreground=[('selected', COLORS['accent_blue'])]
        )

    def create_styled_button(
        self, parent, text, command, button_type="blue", **kwargs
    ):
        """Create a flat, modern button with hover micro-animations."""
        type_colors = {
            'blue': (COLORS['accent_blue'], COLORS['bg_primary'], 
                     COLORS['accent_hover']),
            'green': (COLORS['accent_green'], COLORS['bg_primary'], 
                      '#0ee58e'),
            'red': (COLORS['accent_red'], COLORS['text_primary'], 
                    '#ff3546'),
            'purple': (COLORS['accent_purple'], COLORS['text_primary'], 
                       '#791dcf')
        }
        
        bg_color, fg_color, hover_color = type_colors.get(
            button_type, type_colors['blue']
        )
        
        btn = tk.Button(
            parent, text=text, command=command, bg=bg_color, fg=fg_color,
            font=("Segoe UI", 10, "bold"), relief=tk.FLAT, bd=0, cursor="hand2"
        )
        
        def on_enter(e):
            if btn['state'] != tk.DISABLED:
                btn.config(bg=hover_color)
                
        def on_leave(e):
            if btn['state'] != tk.DISABLED:
                btn.config(bg=bg_color)
                
        btn.bind("<Enter>", on_enter)
        btn.bind("<Leave>", on_leave)
        
        return btn
\end{lstlisting}

\end{document}
